% Inclusion fragment. Requires amsmath, amssymb, and amsthm.
% New scalar-only completion of the UR equality-reflection bootstrap.
% No equation induction, pairing, general PR compiler, or equality test
% is assumed in this fragment. The predecessor bootstrap is a prerequisite.

\section{Actual equality from scalar iteration uniqueness}
\label{sec:ur-reflection-completion}

Work in intuitionistic first-order logic with actual equality, the scalar
iterator computation equations, numerical zero/successor coverage, and
$0\ne S0$. Assume the contextual scalar uniqueness rule
\[
 \frac{\Gamma\vdash f(Sz)=h(f(z))}
      {\Gamma\vdash f(t)=I(t,f(0);u.h(u))},
 \tag{UR}
\]
where $z$ is fresh for $\Gamma,t$ and all passive parameters. The update
$h$ has no active counter argument. Every displayed function in what
follows is a fixed finite scalar term abbreviation; there are no function
variables and no hidden product-state iteration.

The preceding scalar predecessor construction supplies a term $P$ with
$P0=0$ and $P(Sx)=x$. Addition, defined by
$a+b=I(b,a;u.Su)$, satisfies the commutative-monoid equations by common-update
uniqueness. Define
\[
 d(a,b)=I(b,a;u.Pu),\qquad
 Z(a)=C(a,1,0),\qquad \operatorname{sg}(a)=C(a,0,1),
 \qquad C(v,a,b)=I(v,a;u.b).
\]
The following elementary equations use only UR and computation:
\begin{equation}
\begin{aligned}
 d(a,0)&=a,&d(a,Sb)&=P(d(a,b)),\\
 d(Sa,Sb)&=d(a,b),&d(a,a)&=0,&d(0,a)&=0.
\end{aligned}
 \label{eq:ur-monus-basics}
\end{equation}
For cancellation, the two terms $b\mapsto d(Sa,Sb)$ and $b\mapsto d(a,b)$
have common initial value $a$ and common update $P$. The diagonal
$a\mapsto d(a,a)$ has the identity update and initial value zero. The
constant-zero function and $a\mapsto d(0,a)$ have common update $P$ and
initial value zero. Thus no equation-induction assumption is hidden here.

\subsection{A diagonal descent principle for numerical zero tests}

UR implies the following constant-recurrence principle:
\[
 \Gamma\vdash f(Sz)=f(z)\quad\Longrightarrow\quad
 \Gamma\vdash f(t)=f(0).
 \tag{Const}
\]
Apply UR with identity update to $f$ and to the constant function $f(0)$.

Fix a finite arity $r\geq1$ and a scalar term $q(x_1,\ldots,x_r)$.
Suppose that, with any passive parameters fixed,
\[
 q(Sx_1,\ldots,Sx_r)=q(x_1,\ldots,x_r)
\]
and that $q$ is zero whenever any of its displayed arguments is zero.
Then
\[
                         q(x_1,\ldots,x_r)=0.             \tag{DD}
\]
Indeed, coverage in each argument gives
$q(Px_1,\ldots,Px_r)=q(x_1,\ldots,x_r)$: if any argument is zero,
both sides are zero; otherwise the successor identity applies.
Now put
\[
 f(k)=q(d(x_1,k),\ldots,d(x_r,k)).
\]
The defining equation for $d$ gives $f(Sk)=f(k)$, so Const makes $f$
constant. At $k=x_1$ its first argument is $d(x_1,x_1)=0$; hence
$f(x_1)=0$. At $k=0$ it is the desired $q(x_1,\ldots,x_r)$.
This argument uses scalar common-update uniqueness. It does not quantify over
tuples or invoke a simultaneous-recursion rule.

For Boolean-valued terms $p,q$, use
\[
 \operatorname{or}(p,q)=C(p,q,1),\qquad
 \operatorname{xor}(p,q)=C(p,q,Z(q)).
\]
Here Boolean-valued means that the term is provably either $0$ or $1$.
Finite cases show that $\operatorname{xor}(p,q)=0$ if and only if $p=q$.
This finite Boolean fact is not reflection of arbitrary numerical equality.

\subsection{Two Boolean comparison identities}

First, DD gives
\[
 Z(d(Sn,a))=\operatorname{sg}(d(a,n)).
 \tag{B1}\label{eq:ur-strict-complement}
\]
To verify its DD premises, let $q(n,a)$ be the Boolean xor of the two
displayed sides. Simultaneously replacing $n,a$ by $Sn,Sa$ cancels one
successor in each subtraction, so $q(Sn,Sa)=q(n,a)$. If $a=0$, both
sides are zero. If $n=0$, split $a$ into zero and successor: the former
case again gives two zeros, while $a=Sb$ gives
$d(1,Sb)=d(0,b)=0$ and $\operatorname{sg}(Sb)=1$, so both sides are one.
Thus both zero-coordinate boundaries vanish and DD applies.

Second, DD gives
\[
 Z\bigl(d(d(x,a),d(b,a))\bigr)
 =\operatorname{or}\bigl(Z(d(x,a)),Z(d(x,b))\bigr).
 \tag{B2}\label{eq:ur-comparison-join}
\]
Again use the Boolean xor $q(x,a,b)$ of the two sides.
When all three arguments are successors, simultaneous cancellation
reduces the expression to $q(x,a,b)$.
The three zero-coordinate boundaries are as follows.
If $x=0$, both displayed sides are one.
If $a=0$, the left side is $Z(d(x,b))$, and the right side is
$\operatorname{or}(Z(x),Z(d(x,b)))$; coverage on $x$ makes them equal.
If $b=0$, the left side is $Z(d(x,a))$, and the right side is
$\operatorname{or}(Z(d(x,a)),Z(x))$; coverage on $x$ again proves equality.
In both latter arguments, the $x=0$ branch makes the other Boolean one,
whereas the successor branch makes $Z(x)=0$.
Consequently DD proves the vanishing of the xor, and Boolean cases
give the stated identity.

Coverage and the successor equation for addition also give
\[
 \operatorname{or}(\operatorname{sg}(u),\operatorname{sg}(v))
       =\operatorname{sg}(u+v).
 \tag{B3}\label{eq:ur-sign-or}
\]
For $u=0$ this reduces to the sign of $v$; for successor $u$, both
sides are one.

\subsection{An explicit scalar accumulator}

Define the two scalar terms
\begin{equation}
\begin{aligned}
 \mu(n,a)&=d(n,d(n,a)),\\
 \phi(n,a,b)&=\mu(n,a)+\mu(d(n,a),d(b,a)).
\end{aligned}
 \label{eq:ur-scalar-accumulator}
\end{equation}
These are finite compositions of $d$ and addition. In particular,
$\phi$ is not introduced by a counter-dependent recursion.
We now prove its useful counter-dependent equation from its explicit term.

Coverage on $a$ and the monus equations give
\[
                         P(d(Sn,a))=d(n,a).
 \tag{Pd}
\]
At $a=0$ this is $P(Sn)=n$. At $a=Sa'$ it follows from
$d(Sn,Sa')=d(n,a')$ and $P(d(n,a'))=d(n,Sa')$.
Split the value $r'=d(Sn,a)$ into zero and successor.
If $r'=0$, then $d(n,a)=0$ by Pd, so
$\mu(Sn,a)=Sn$ and $\mu(n,a)=n$.
If $r'=Sr$, then $d(n,a)=r$ and
$\mu(Sn,a)=d(Sn,Sr)=d(n,r)=\mu(n,a)$.
In both cases,
\[
             \mu(Sn,a)=\mu(n,a)+Z(d(Sn,a)).              \tag{M}
\]

Put $h=d(b,a)$ and repeat the same split of $r'=d(Sn,a)$.
If $r'=0$, the first argument of the second $\mu$ term is zero at both
$n$ and $Sn$, and $\phi(Sn,a,b)=Sn$, $\phi(n,a,b)=n$.
If $r'=Sr$, the first summand is unchanged, while M gives
\[
                 \mu(Sr,h)=\mu(r,h)+Z(d(Sr,h)).
\]
Associativity therefore proves, in both cases,
\[
 \phi(Sn,a,b)=\phi(n,a,b)+Z(d(d(Sn,a),d(b,a))).
\]
Apply B2 at $x=Sn$, then B1 to each Boolean on its right, and finally B3.
The resulting recursion specification is
\[
 \boxed{\quad
 \phi(0,a,b)=0,\qquad
 \phi(Sn,a,b)=\phi(n,a,b)+\operatorname{sg}(d(a,n)+d(b,n)).
 \quad}                                                   \tag{Acc}
\]
All identities leading to Acc used scalar uniqueness, coverage, and
finite case analysis. No equality-reflection or induction theorem has
yet been used.

\subsection{Commutativity of the arithmetic maximum expression}

Put
\[
                              m(a,b)=a+d(b,a).
\]
The name is descriptive only; its order-theoretic properties are not
assumptions. Four zero/successor cases prove
\[
                  m(a,b)=m(Pa,Pb)+\operatorname{sg}(a+b).
 \tag{Red}\label{eq:ur-max-diagonal-reduction}
\]
At $(0,0)$ both sides are zero. At $(0,Sb)$ and $(Sa,0)$ it is the
successor equation for addition. At $(Sa,Sb)$, simultaneous subtraction
cancellation gives
$m(Sa,Sb)=Sa+d(b,a)=S(m(a,b))$, and the sign of the sum is one.

For fixed $a,b$, form the scalar terms
\begin{align*}
 F(n)&=m(d(a,n),d(b,n))+\phi(n,a,b),\\
 G(n)&=m(d(b,n),d(a,n))+\phi(n,a,b).
\end{align*}
Red, Acc, and commutativity of addition give $F(Sn)=F(n)$ and
$G(Sn)=G(n)$. Hence Const gives
\[
                          F(n)=m(a,b),\qquad G(n)=m(b,a).
\]
At $n=b$, the equality $d(b,b)=0$ and the identities
$m(c,0)=c=m(0,c)$ make both terms equal, by these equations, to
$d(a,b)+\phi(b,a,b)$. It follows that
\[
                    \boxed{\ a+d(b,a)=b+d(a,b).\ }       \tag{Max}
\]
This is the common-accumulator argument for maximum commutativity.
The scalar-specific ingredient is the explicit accumulator in
\eqref{eq:ur-scalar-accumulator} and its derivation through Boolean DD;
no prior availability of arbitrary primitive recursion is required.

\subsection{Reflection and recovery of equation induction}

Define
\[
                       \Delta(a,b)=\operatorname{sg}(d(a,b)+d(b,a)).
\]
Coverage and separation imply $\operatorname{sg}(u)=0\leftrightarrow u=0$.
They also imply $u+v=0\leftrightarrow(u=0\land v=0)$: a successor
in either summand would make the sum a successor, contradicting separation.
Thus $\Delta(a,b)=0$ forces $d(a,b)=0=d(b,a)$. Max then yields
\[
               a=a+d(b,a)=b+d(a,b)=b.
\]
Conversely $a=b$ gives both subtractions zero. Therefore
\[
 \boxed{\quad \Delta(a,b)=0\ \longleftrightarrow\ a=b,\qquad
                  \Delta(a,b)=0\ \lor\ \Delta(a,b)=1.\quad}
 \tag{EqRef}
\]
The equality on the right of EqRef is the logic's actual equality on
the whole numerical domain. Nothing restricts $a,b$ to standard numerals
or to the image of an encoding.

Finally, the independently proved UR retraction lemma gives contextual
induction for all zero equations $q(n)=0$. Given an induction premise
for any equation $s(n)=t(n)$, use EqRef to transfer its base and step
to the zero equation $\Delta(s(n),t(n))=0$, apply that zero-induction
rule, and transfer the conclusion back. This recovers contextual
literal-equation induction in arbitrary retained first-order contexts.
The reverse derivation of UR from literal-equation induction is immediate
by induction on $f(n)=I(n,f(0);h)$. Consequently scalar UR and contextual
equation induction are equivalent over the common scalar computation,
coverage, and separation background, once the preceding UR predecessor
bootstrap is included.
